\section{Experimental Setup}
\label{sec:setup}

\paragraph{Benchmarks.} We evaluate on four benchmarks whose answers depend on images. MemEye~\citep{guo2026memeye} contains 371 questions over eight scenario histories, with 126.5 rounds per history on average. MemLens~\citep{ren2026memlens} gives each question its own history, averaging 64.6 rounds at 32k and 124.8 rounds at 64k. Its gold rounds belong to the answer sessions and contain an evidence turn or image. We report its 466 test questions; the other 233 serve as development questions. VisualMem~\citep{nguyen2026personalvisual} contains 696 multiple-choice questions across ten personas, whose histories average 1,016.6 rounds and 172 images. The questions concern facts shown only in images and primarily test whether a method can locate one answer-bearing image in a long history. CamRoll~\citep{nguyen2026camroll} contains 20 personal camera rolls of real photos from YFCC100M, 780 photos each on average, with 1,000 multiple-choice questions about the album owner and about specific photos. Each photo is a round whose text is its capture time, and every reader also receives the owner's reference photo. For both multiple-choice benchmarks we score the option letter and downscale images to 1,024 pixels. Prompts and the gate threshold were chosen once on development data (MemEye and the 233 MemLens development questions) and held fixed; the MemLens test questions, VisualMem, and CamRoll test the same configuration on unseen data.

\paragraph{Models.} The lightweight reader is \texttt{gemini-2.5-flash-lite} (temperature 0, no thinking), and the strong reader, the selector, and the captioner are \texttt{gemini-3.8-flash} (low thinking). The note writer is the lightweight model, and \texttt{gemini-embedding-2} embeds each round as its text plus its image descriptions.

\paragraph{Baselines and controls.} RAG@$k$ gives the reader the $k$ most similar rounds with their original images, ordered chronologically; we report $k{=}40$ and give smaller budgets in Appendix~\ref{app:retrieval}. As a stronger retrieval baseline, the selector model re-ranks the 100 most similar rounds with the selector instruction, and the reader receives the top 40 (\emph{LLM rerank}; Table~\ref{tab:rerank}). Full-context reading supplies the whole history. Selection controls use the same selector and the same 100-round candidate pool as \qss{}: the re-ranking cut at eight rounds or at the size that \qss{} chose, and cosine retrieval of that size.

\paragraph{Memory systems.} We compare with five recent systems. SimpleMem~\citep{liu2026simplemem}, Omni-SimpleMem~\citep{liu2026omnisimplemem}, and TeleMem~\citep{chen2026telemem} write text records at ingestion; V-Mem~\citep{kang2026vmem} and MMA~\citep{lu2026mma} keep raw turns and retrieve them with modality routing and confidence weighting. All systems build memory from the same rounds and image descriptions and send their outputs to the same readers, prompts, and judges (Appendix~\ref{app:baselines}).

\paragraph{Metrics and statistics.} We measure retrieval by full-set recall (Eq.~\ref{eq:fsr}) and the number of opened rounds. Answer accuracy follows the official benchmark rubrics, using a Gemini judge for MemEye's 0--1 scores and Qwen3-VL-235B for MemLens's binary scores. A second judge family preserves the main method comparisons on MemEye (Appendix~\ref{app:robust}). We test paired differences over questions with a sign-flip permutation test and give 95\% bootstrap intervals in the appendix. Costs use measured tokens at list prices, per 1,000 questions. Appendix~\ref{app:latency} reports token F1, BLEU-1, Recall@10, nDCG@10, and latency.
